% Fichier généré automatiquement par tools/generate_corrections.py.
% Source : CIN/CIN-02-VitesseAcceleration/06_TR_02/06_TR_02.tex
% Statut : complet (3/3 question(s) renseignée(s), 0 reprise(s)).
\begin{corrigebox}[Corrigé extrait de la banque]
\CorrigeQuestion{1}
\textbf{Méthode 1 -- Dérivation vectorielle} ~\\

$\vectv{C}{2}{0} =$
$\deriv{\vect{AC}}{\rep{0}}= \deriv{\vect{AB}}{\rep{0}}+\deriv{\vect{BC}}{\rep{0}}$
$= \dot{\lambda}(t) \vect{i_0} + R \deriv{\vi{2}}{\rep{0}}$
$= \dot{\lambda}(t) \vect{i_0} + R \dot{\theta} \vj{2}$

\textbf{Méthode 2 -- Composition du torseur cinématique} ~\\
$\vectv{C}{2}{0} =\vectv{C}{2}{1} +\vectv{C}{1}{0} $

Pour tout point $P$, $\vectv{P}{1}{0}=\dot{\lambda}\vi{0}$.

$\babarv{C}{B}{2}{1}=-R\vi{2}\wedge \dot{\theta}\vk{0}=R\dot{\theta}\vj{2}$.

On a donc $\vectv{C}{2}{0}=\dot{\lambda}\vi{0}+R\dot{\theta}\vj{2}$.
\CorrigeQuestion{2}
$\torseurcin{V}{2}{0}=\torseurl{\vecto{2}{0}=\dot{\theta}\vk{0}}{\vectv{C}{2}{0}=\dot{\lambda}\vi{0}+R\dot{\theta}\vj{2}}{C}$.
\CorrigeQuestion{3}
$\vectg{C}{2}{0} =\deriv{\vectv{C}{2}{0}}{\rep{0}}$
$= \ddot{\lambda}(t) \vect{i_0} + R \deriv{\dot{\theta} \vj{2}}{\rep0}$
$= \ddot{\lambda}(t) \vect{i_0} + R \left(\ddot{\theta} \vj{2} - \dot{\theta}^2 \vi{2}\right) $.
\end{corrigebox}
